\chapter{العلل في الأحاديث واختلاف الروايات}

\section{مقدمة في اختلاف الأحاديث وموافقة السنة للقرآن}
الحمد لله والصلاة والسلام على رسول الله، أما بعد؛ فقد عقد الإمام الشافعي رحمه الله تعالى باباً عظيماً أسماه «باب العلل في الأحاديث»، والعلة في اصطلاح المحدثين: سبب خفي قادح يقدح في صحة الحديث مع أن ظاهره السلامة، والتقييم هنا أعم من العلة الاصطلاحية، فهو يقصد أسباب الاختلاف الظاهري بين الأحاديث.

وقد سأل سائل الإمام الشافعي عن أحاديث ظاهرها الاختلاف والتضاد؛ منها ما يبين القرآن، ومنها ما يزيد عليه، ومنها ما ينهى عنه النبي صلى الله عليه وسلم فيعده العلماء تحريماً وفي موضع آخر يعدونه كراهة، ومنها ما يؤخذ به ويترك ما يعارضه. فما هي الحجة والقاعدة في هذا الأخذ والترك والترجيح؟

أجاب الشافعي بتأصيل عظيم: «كل ما سن رسول الله صلى الله عليه وسلم مع كتاب الله، فهو موافق لكتاب الله». فإما أن تأتي السنة بنص يوافق نص القرآن، وإما أن تبين مجمل القرآن وتفصله، وإما أن تأتي بسنة ليس فيها نص في القرآن. وفي كل هذه الحالات، يجب التسليم لأن الله هو الذي فرض علينا طاعة رسوله صلى الله عليه وسلم.

\section{الناسخ والمنسوخ في السنة}
كما أن القرآن ينسخ بعضه بعضاً (كنسخ المصابرة من واحد لعشرة إلى واحد لاثنين)، كذلك سنة رسول الله صلى الله عليه وسلم ينسخ بعضها بعضاً. ومثاله قوله صلى الله عليه وسلم: «كنت نهيتكم عن زيارة القبور فزوروها»، وقوله: «كنت نهيتكم عن ادخار لحوم الأضاحي بعد ثلاث فكلوا وادخروا». فالسنة ينسخ بعضها بعضاً لأن الكل وحي من عند الله عز وجل.

\section{أسباب الاختلاف الظاهري بين الأحاديث}
لا يوجد في أحاديث رسول الله صلى الله عليه وسلم الصحيحة اختلاف تضاد أبداً، وإنما يقع الاختلاف في فهم السامع لعدة أسباب بينها الشافعي:

\subsection{العام المراد به الخصوص}
رسول الله صلى الله عليه وسلم عربي اللسان والدار، وقد يقول القول عاماً ويريد به العام، وقد يطلق اللفظ عاماً ويريد به الخاص، فمن لم يفهم لغة العرب توهم التعارض.

\subsection{سماع الجواب دون السؤال}
قد يُسأل النبي صلى الله عليه وسلم عن شيء فيجيب على قدر المسألة، فيحضر الراوي فيسمع الجواب ولا يدرك السؤال ومناطه، فيحدث بالجواب مطلقاً. كمن روى قوله: «هو الطهور ماؤه الحل ميتته» دون أن يذكر أنه كان جواباً لمن سأل عن الوضوء بماء البحر لتوفير الماء العذب للشرب.

\subsection{اختلاف الحالين}
قد يسن النبي صلى الله عليه وسلم في شيء حكماً، وفيما يخالفه حكماً آخر، فلا يفرق بعض السامعين بين اختلاف الحالين اللتين سن فيهما. 
ومثاله بيوع العرايا؛ فقد نهى النبي صلى الله عليه وسلم عن الغرر والجهالة والمزابنة، ولكنه رخص لبعض الفقراء الذين لا يملكون نقوداً أن يشتروا الرطب على النخل بخرصه من التمر ليأكلوه (بيوع العرايا). فمن لم يفرق بين الحالين ظن التعارض، وإنما العرايا حالة رخصة مخصوصة لحاجة، والنهي العام باقٍ على التحريم.

\subsection{حفظ المنسوخ دون الناسخ}
قد يحفظ بعض الصحابة الحديث المنسوخ ويغيب عنه الناسخ، فيحدث بما حفظ. ويحفظ صحابي آخر الناسخ فيحدث به. وهذا لا يذهب على عامة الصحابة، فلا تضيع سنة ناسخة عن مجموع الأمة أبداً، فإذا جمعت طرق الحديث تبين الناسخ من المنسوخ.

\section{وجوب التسليم للسنة وترك الاعتراض}
يقرر الشافعي قاعدة إيمانية جليلة: «لا يجوز أن تقول لسنة رسول الله ﷺ: لِمَ ولا كيف؟». 
فلا يجوز لمسلم أن يفرق بين ما جمعه النبي صلى الله عليه وسلم، ولا أن يجمع بين ما فرقه. ومن اعترض على حكم رسول الله صلى الله عليه وسلم متسائلاً على وجه الإنكار "لماذا فرق بين هذا وهذا؟" فهو إما جاهل، أو مرتاب شاك، والشك والارتياب في كلام النبي صلى الله عليه وسلم شر من الجهل وقد يؤدي للكفر.
فالمؤمن يقول: «سمعنا وأطعنا»، ويسلم لحكم الله ورسوله دون حرج. أما ادعاء بعض المنتكسين في أيامنا أنهم "سلفيون عقلانيون" يقدمون العقل على النقل، فهو طعن في السلفية وبوابة للعلمانية، فالعقل الصريح لا يعارض النقل الصحيح أبداً.

\section{طرق دفع التعارض (الترجيح)}
إذا وجدنا حديثين ظاهرهما الاختلاف، فإننا نجمع بينهما بالوجوه السابقة، فإن لم يمكن، صرنا إلى الترجيح بناءً على أسس منها:
\begin{itemize}
    \item أن يكون أحد الحديثين له دلالة ومواكبة من كتاب الله أو سنة نبيه تدعمه.
    \item أن يكون أحد الحديثين يوافق القواعد الكلية والأصول الثابتة.
    \item أن يكون رواته أحفظ وأكثر ضبطاً.
\end{itemize}
والأصل أن ما نهى عنه رسول الله صلى الله عليه وسلم فهو على التحريم، حتى تأتي دلالة تصرفه إلى الكراهة أو التنزيه.